\chapter{Lire, chercher et démontrer}
\label{exercices:01}
\section*{Exercice 1.1 -- Vrai, faux, ou données insuffisantes}
Pour chaque affirmation, donner une preuve ou un contre-exemple.
\begin{enumerate}\item Tout multiple de $6$ est multiple de $3$.
\item Tout multiple de $3$ est multiple de $6$.
\item Si deux triangles ont les mêmes angles, ils ont les mêmes côtés.
\item Pour tout entier $n$, $n(n+1)$ est pair.\end{enumerate}

\section*{Exercice 1.2 -- La bonne implication}
Écrire la réciproque de : « Si un quadrilatère est un carré, alors ses quatre côtés sont égaux. » Est-elle vraie ? Quel mot décrit un quadrilatère à quatre côtés égaux ?

\section*{Exercice 1.3 -- Démontrer avec une lettre}
Démontrer que la somme de deux entiers impairs est paire. Expliquer pourquoi vérifier $3+7=10$ ne suffit pas.

\section*{Exercice 1.4 -- Réparer une rédaction}
Un élève écrit : « Sur ma figure, le triangle a l'air rectangle. J'utilise donc Pythagore. » Que manque-t-il ? Proposer deux types de données qui permettraient de poursuivre correctement.
